\begin{abstractCN}
\setlength{\baselineskip}{23pt} %设置行距23磅

随着物联网的发展，状态更新业务不断涌现，对这类业务而言，接收端更关注信息的新鲜度。然而，吞吐量和时延却难以对其进行衡量，为解决这个问题，信息年龄（Age of Information，AoI）被提出并受到广泛关注。随着物联网部署规模的扩张，状态更新业务的设备数量也迅速增长。在这种场景下，集中式调度需要维护大量节点状态，控制和反馈的开销较高。相比之下，随机接入无需掌握全局状态，更适合海量设备动态接入的场景。正因如此，如何在大规模接入的条件下兼顾接入效率与信息新鲜度，是一个值得研究的 MAC 层问题。

然而，现有研究没有充分考虑大规模接入场景下持续碰撞所带来的影响，而且大多都将重点放在平均 AoI，对 AoI 违规概率的关注较少。对于工业控制等时间敏感型业务，仅考察平均 AoI 不足以判断系统的可靠性。针对上述问题，本文充分考虑了大规模接入场景下持续碰撞所带来的负面影响，先是设计了一种基于等待的随机接入退避机制，并在此基础上，围绕时隙 ALOHA 和标准 CSMA 两类随机接入框架来构建相应的理论分析模型，推导出平均 AoI 与 AoI 违规概率的闭式表达式。本文的主要工作如下：

\begin{enumerate} [wide]
    \item \textbf{设计面向大规模接入场景的等待机制。}针对现有研究对持续碰撞关注不足的问题，本文设计了一种基于等待的随机接入退避机制，通过规划节点在成功更新以及遭遇碰撞后的不同退避行为，缓解信道拥塞。
    \item \textbf{将所提机制引入时隙 ALOHA 框架中并进行性能分析。} 针对节点状态耦合的难题，本文通过解耦的近似方法，将高维复杂的全局状态降维成单节点的马尔可夫链模型，在此基础上建立并求解系统的稳态方程。随后，基于求得的稳态特征，推导出平均 AoI 及 AoI 违规概率的闭式表达式，并分析系统参数对性能的影响。
    \item \textbf{将所提机制引入标准 CSMA 框架中并进行性能分析。} 针对节点状态耦合，以及由于载波侦听、退避冻结导致的系统时隙变长两个难题，本文通过解耦的近似方法，构建单节点的马尔可夫链模型，在此基础上建立并求解系统的稳态方程。随后，将稳态特征与变长时隙的统计特性相结合，推导出平均 AoI 及 AoI 违规概率的闭式表达式，并分析系统参数对性能的影响。
    \item \textbf{验证理论模型并揭示等待机制的性能演化规律。}本文通过仿真验证了理论表达式的准确性，并与多种随机接入机制进行了对比。结果表明，在中高负载尤其是海量设备接入的场景下，适度等待并不会损害信息新鲜度。相反，它能够减少信道中的竞争，从而同时改善平均 AoI 与 AoI 违规概率。
\end{enumerate}

本文的研究不仅丰富了面向信息年龄的随机接入理论分析体系，也为未来状态更新系统的协议设计与优化提供了有价值的理论参考与工程指导。

\end{abstractCN}


\begin{keywordCN}
随机接入，信息年龄，性能分析
\end{keywordCN}

\begin{abstractEN}
\setlength{\baselineskip}{23pt} %设置行距23磅

With the development of the Internet of Things (IoT), status update services continue to emerge. For these services, the receiver focuses more on the freshness of the information. However, traditional metrics like throughput and delay are difficult to measure information freshness. To solve this problem, a new metric called Age of Information (AoI) was proposed and has received widespread attention.With the expansion of IoT deployment, the number of devices for status update services is growing rapidly. In such scenarios, centralized scheduling requires maintaining the states of a massive number of nodes, which incurs high control and feedback overhead. In contrast, random access eliminates the need for global state information, making it more suitable for the dynamic access of massive devices. Therefore, how to balance access efficiency and information freshness under large-scale access conditions is a crucial MAC layer problem.

However, existing research does not fully consider the impact of continuous collisions in large-scale access scenarios, and most studies focus on the average AoI, paying less attention to the AoI violation probability. For time-sensitive services such as industrial control, considering only the average AoI is not enough to determine system reliability. To address these problems, this thesis fully considers the negative impact of continuous collisions in large-scale access scenarios. First, we design a waiting-based random access backoff mechanism. Based on this, we build corresponding theoretical analysis models for two random access frameworks: slotted ALOHA and standard CSMA, and derive closed-form expressions for the average AoI and the AoI violation probability. The main contributions of this thesis are as follows:

\begin{enumerate} [wide]
\item \textbf{Designing a waiting mechanism for large-scale access scenarios.} To address the problem that existing research pays insufficient attention to continuous collisions, we design a waiting-based random access backoff mechanism. By planning different backoff behaviors for nodes after successful updates and collisions, this mechanism effectively relieves channel congestion.

\item \textbf{Introducing the proposed mechanism into the slotted ALOHA framework and analyzing its performance.} To solve the challenge of coupled node states, we use a decoupling approximation method to simplify the complex system into a single-node Markov chain model. Based on this, we establish and solve steady-state equations of the system. Based on the obtained steady-state features, we derive closed-form expressions for the average AoI and AoI violation probability, and analyze the impact of system parameters on performance.

\item \textbf{Introducing the proposed mechanism into the standard CSMA framework and analyzing its performance.} To address the challenges of coupled node states and variable-length slots caused by carrier sensing and backoff freezing, we use a decoupling approximation method to build a single-node Markov chain model. Based on this, we establish and solve the steady-state equations of the system. Then, by combining the steady-state features with the statistical characteristics of the variable-length slots, we derive closed-form expressions for the average AoI and the AoI violation probability, and analyze the impact of system parameters on performance.

\item \textbf{Verifying the theoretical model and revealing the performance evolution rules of the waiting mechanism.} Through simulation experiments, we verify the accuracy of the theoretical expressions. We also compare our mechanism with several random access mechanisms. The results show that under medium-to-high loads, especially in massive device access scenarios, moderate waiting does not harm information freshness. Instead, it reduces channel competition, thereby improving both the average AoI and the AoI violation probability simultaneously.
\end{enumerate}

In summary, our research not only enriches the theoretical analysis of random access for Age of Information, but also provides valuable theoretical reference and engineering guidance for the protocol design and optimization of future status update systems.

\end{abstractEN}

\begin{keywordEN}
Random access, Age of information, Performance analysis
\end{keywordEN}